\section{Retained-relation operativity and its applications}
\label{manual:m2-awareness-fixtures}
\subsection{Exact model and matched state}
\label{sec:m2-awareness-ladder}
These witnesses instantiate the Main state, retention and matched causal
criterion in \mainequation{main:eq:awareness-state-factorization},
\mainequation{main:eq:retained-rebalancing} and
\mainequation{main:eq:awareness-active}. They are operational finite models. All inputs and
expected results are in
\texttt{manual/data/m2/awareness\_cases.jsonl}; the shared runner is
\texttt{manual/scripts/aod\_awareness.py}. From the source root, run
\begin{quote}\small\ttfamily
python3 manual/scripts/aod\_awareness.py --check-data
\end{quote}
Actual state and outcome traces are separate execution evidence. Expected
answers, row names and passive observations are never arguments of the
response law.

The model instantiates \mainequation{main:eq:awareness-state-factorization}.
Its present-record memory and retention placeholders are empty; \(M^C\) stores
return identities and \(R\) alone stores operative prior rebalancing.
The world validates admitted waves, issued actions and prior-operand commitments;
it cannot recreate a removed residual for the update. No passive record or
cache may restore a removed value; retained bytes, fixture verdicts and role
labels are excluded from the outcome. Prior provenance also binds any action
witness. An operation-count proxy records work, not operativity.

\subsubsection{Rational response and continuation}
\paragraph{Setup.}
Consider the following rational response realization on already admitted
support, with \(p>0\), gain \(g\), available source set \(D_k\), and explicitly
perturbative query \(q_k\):
\begin{equation}
 J_{k+1}=J_k+\frac{A_k+q_k+g\sum_{e:\,i_e\in D_k}r_e}{p}.
 \label{eq:awareness-response}
\end{equation}
The term \(J_k+A_k/p\) specializes \mainequation{main:eq:sadar-rational-response};
\(A_k\) is a scalar attention component, separate from \(r_e\).
Its channel construction and \(AO=-OA\) remain fixed.

\paragraph{Derivation.}
Subtract the two matched
responses on common source/event support, extending absent summands by zero:
\begin{equation}
 J_{k+1}(R)-J_{k+1}(R')
 =\frac{g}{p}\left(\sum_{e:\,i_e\in D_k}r_e-
                         \sum_{e:\,i_e\in D_k}r'_e\right).
 \label{eq:awareness-causal-difference}
\end{equation}
A nonzero right side therefore proves a later current difference for this
model. At \(J_k=A_k=q_k=0\), \(g=1\), \(r=1\), the retained and removed cases
give \((1,0)\) for \(p=1\), and \((1/2,0)\) for \(p=2\). These are two
separate pressure families, each with its own matched retention comparison.
Zero gain or zero effective retention gives no such witness.

\paragraph{Specialization.}
Let present state include a Q4 vertex \(v\) and fibre role
\(\epsilon\in\{-1,0,1\}\). Two candidate continuations go to \(v\oplus1\)
and \(v\oplus2\). A branch enters the hinge; from the hinge the two candidates
carry the opposite branch choices. Both obey the inherited H1 and fibre
rules. A declared sign selector chooses the first candidate when
\(J_{k+1}>0\), otherwise the second. The selector consumes an already
admitted finite support; it does not create H1 or certify closure. Every
such continuation contributes zero new closure-phase cycles until its separate
retained-start return is certified.

\paragraph{Falsification.}
A difference in current pressure, law or future exogenous input invalidates
a retention-only comparison. Recounting \(R\) itself as the outcome,
recovering a removed value through another channel, using an unbound
phase type, or allowing a same-cut effect also fails the criterion.
A label attached to either execution cannot change \(\mathcal U\).

\aodprovenance{\mainequation{main:eq:hinge-grammar};
\mainequation{main:eq:atomic-h1};
\mainequation{main:eq:wave-update-trace};
\mainequation{main:eq:causal-update};
\eqref{eq:awareness-response};
\eqref{eq:awareness-causal-difference}.}

\paragraph{Response executions.}
The target has four completed returns separated by ten native indices. The
ordinary peer has separations \((7,7,7)\). Repeated native contacts bind both
prefixes and the signed count-lag comparison. Exact independent and cross
contributions \((2)\) and \((1)\) yield the retained FRG residual \(r=1\).
The prior contact cut is \(23\), retention follows at \(24\), the complete
matched boundary is \(33\), and the first changed response is at \(34\).
These are native commitment indices, not time measurements.

\begin{center}
\begin{tabular}{lrrrl}
\toprule Pressure & Retained \(r\) & First current & Removed current & Restoration\\\midrule
\(1\) & \(1\) & \(1\) & \(0\) & Exact state and outcome recovery\\
\(2\) & \(1\) & \(1/2\) & \(0\) & Exact state and outcome recovery\\
\(1\), gain \(0\) & \(1\) & \(0\) & \(0\) & No causal witness\\
\(1\) & \(0\) & \(0\) & \(0\) & No causal witness\\
\bottomrule
\end{tabular}
\end{center}
Each row applies \mainequation{main:eq:awareness-active} in its pressure family;
changing pressure across a pair is rejected. A passive trace of the
same execution leaves all state and event bytes unchanged. Ordinary attention
\(A=1\) with empty \(R\) changes current but supplies no retention effect.
A zero-gain retained residual likewise fails the matched operativity criterion.

\subsection{Hostile state, lineage and timing controls}
\label{sec:m2-hostile}
\label{man:temporal-awareness-controls}
The hostile rows reject changed present pressure, closure memory, update law
or future input; duplicate retention inside \(X\); unknown lineage; a
future or noninteger commitment; and a retained value detached from its
prior operands. The executable controls additionally remove motif
certificates, fabricate cadence or phase count, reverse carried orientation,
change the return readout, and forge an issued action. Recommitting altered
bytes in the diagnostic registry does not bypass the structural type and
causal checks.

A relational-phase control uses three actual source lineages. Permuting the
three retained lag values preserves their marginal multiset but destroys the
source-specific native comparisons. Tau, normalized-circle and old selector
coordinates are rejected as substitutions for
\texttt{M1:unwrapped-count-lag}. The comparator has exactly two declared bias
operands. No phase label or scalar memory value supplies admission.
Every selected response is an H1 edge with a lawful fibre transition.
Uncertified continuations add zero closure-phase cycles, including when their
endpoints happen to coincide.

\subsection{Operative boundary, labels, action and learning}
\label{sec:self-proprio-toys}
\paragraph{Setup.}
The operative boundary is \mainequation{main:eq:operative-self}.
Sensor, register, CPU and vocabulary models reuse this interface.
\paragraph{Derivation.}
Two retained source contributions produce a three-node operative subgraph:
two source interfaces and the target. The target's role is SELF, its active contact
interfaces are BETWEEN, and unconnected nodes are NOT-SELF. Null remains
separately typed. With current outcome \(O=J\), nonzero gain and both source
channels available, consider the existing instance \(J_k=A_k=q_k=0\),
\(p=g=1\). Contributions \(+1\) and \(-1\) cancel in total while each
removal changes later current by \eqref{eq:awareness-causal-difference}: the full result is \(0\), removal of the
\(+1\) channel gives \(-1\), and removal of the \(-1\) channel gives \(+1\).
Both source lineages therefore remain operative. This aggregate-zero state,
a provenance-bound zero-valued retained channel, and structural absence have
different records and digests. Neither zero-valued case forces a fold or a
positive sheddic result. The antisymmetric pressure ledger also retains both channels.
A CPU interface uses a
declared cost of three operations per available source; that resource proxy
alone supplies no operativity verdict.

\subsection{Action return and persistent learning}
\label{subsec:proprioception-learning}
\paragraph{Setup.}
An issued action binds an immutable identity to its issuer, native commitment,
command and channel. Both issuer and channel must already have admitted
recurrent prefixes at issuance. A proprioceptive consequence must return on that exact
action lineage with its own certified closure. An observed value is supplied
by a declared projection of the returned trace, not a free scalar attached
to an unchanged return. In a native count realization, write
\begin{equation}
 a=(\mathrm{id},i,k_a,u,\mathrm{channel}),\qquad
 y=\#\{\text{atomic traversals in the bound return}\},\qquad r_a=y-u.
 \label{eq:action-return-retention}
\end{equation}
The command \(u\) and readout \(y\) share this declared count domain. They are
not measured physical quantities. The action precedes the return, the
retention event follows its committed input, and a later transition consumes
\(r_a\). The action's immutable identity prevents conflicting command records
from masquerading as the same issuance.

\paragraph{Derivation.}
For command \(u=1\) and a certified two-traversal return, \(y=2\) and
\(r_a=1\). The matched response difference is then \(g/p\).
Uncommanded sensing, motor output alone, a return from another action or an
unrelated sensor magnitude cannot establish this proprioceptive relation.
For the single acquired
relation \(r_a=1\), with its source available throughout, \(A_k=q_k=0\) and
\(g=p=1\), the retained model advances current by one per transition, while
the matched removed model retains its current value.
Restoration of the acquired coordinate restores the later increments.

\paragraph{Falsification.}
A surviving record with no later behavioral change is persistence alone.
A changed stimulus cannot be attributed to memory by an unmatched comparison.

\aodprovenance{\mainequation{main:eq:closure-certificate};
\mainequation{main:eq:retained-rebalancing};
\eqref{eq:awareness-causal-difference}.}

\paragraph{Action and learning executions.}
The sensor--motor witness issues its command at \(13\), after the issuer and
channel have acquired recurrent prefixes at \(12\) and \(9\). Its certified
action return occupies \(14,15,16\), with exactly two atomic traversals.
The readout is the declared traversal-count projection \(y=2\), the command
is \(u=1\), and the retained residual is \(y-u=1\). The same certified return
cannot be assigned an arbitrary readout of \(200\). Wrong-action, unissued,
wrong-return and future-return cases fail. Passive sensing and motor output
alone do not establish proprioception.

The separate learning witness uses the already retained relation supplied by
\texttt{core()}, without calling action-return acquisition. It retains that
relation through four zero-stimulus updates, reaching current four. At a new matched boundary, retained execution
then yields currents \((5,6,7)\); removal yields \((4,4,4)\). Exact restoration
of value and lineage restores \((5,6,7)\). The future attention inputs remain
zero in each family, so a changed stimulus does not explain the difference.

\subsection{Polycadence, crescendo, exposure and a second causal link}
\label{sec:m2-crescendo}
The same retention law executes with locked cadences \((10,10,10)\), a declared
nonzero comparator bias \(1/4\) at equal cadence, unequal cadences
\((10,10,10)\) and \((7,7,7)\), or jittered separations \((8,7,13)\).
Each positive retention pair has a later effect. Locked cadence with zero
gain does not. The nonzero bias is an explicitly declared native comparison
coordinate bound to its two operands; it is not circular phase or metrology.

The availability sequence exposes one, two, three, zero, then three source
lineages. The number of operative sources is exactly \((1,2,3,0,3)\), while
all three valid retained records survive dropout. This measures changing
causal participation, not a scalar amplitude or a synchronization score.

\subsubsection{Response-to-load bridge and restoration}
\paragraph{Setup.}
The executable toy separately declares a response-to-load bridge:
\begin{equation}
 L^{\rm toy}_{k+1}=C^{\rm cap}_{k+1}+J_{k+1},\qquad
 C_{\rm close}^{\rm toy}=C^{\rm cap}_{k+1},\qquad
 \Delta_{\rm close}^{\rm toy}=L^{\rm toy}_{k+1}-C_{\rm close}^{\rm toy},
 \quad X_{\rm shedding}^{\rm toy}=\max(0,\Delta_{\rm close}^{\rm toy}).
 \label{eq:awareness-shedding-bridge}
\end{equation}
This model-specific load does not rename \(J\) as SADAR pressure. A positive
record binds source state/event, outer certificate, operands, route, target,
order and digest before exoshedding exposure. Zero means no shedding; negative
residual is deficit/open. Status labels compute neither result.

\paragraph{Derivation.}

With \(O=O_{\rm shed}\) projecting that computed record,
\begin{equation}
 \operatorname{Op}^{O_{\rm shed},H}_e(S_k)=1
 \quad\land\quad
 \operatorname{Restore}_e(S_k^{-e})=S_k
 \label{eq:awareness-conditioned-shedding}
\end{equation}
specializes the causal criterion to computed shedding with exact restoration.
Shedding alone is an M1 consequence; an operative effect without surplus has
no shedding claim.

\paragraph{Falsification.}
A contact, lone outer return or recurring wave cannot replace the fold and
causal witnesses of \mainequation{main:eq:fold-outer-closure},
\eqref{eq:awareness-conditioned-shedding}, and the reapplication of
\mainequation{main:eq:awareness-active} at the outer scope of
\mainequation{main:eq:awareness-carriage}.

\aodprovenance{\mainequation{main:eq:closure-residual};
\mainequation{main:eq:shedding-excess};
\eqref{eq:awareness-shedding-bridge}.}


\paragraph{Fold and shedding executions.}
In the contact/fold witness the first retained response at \(34\) selects a
later outer-construction opportunity. Each outer closure binds the actual
inner certificates and gives one complete closure-phase cycle. Three distinct outer returns permit
recurrence and two distinct higher contacts. The B1 comparison peer also
folds the same admitted inner pair, with three returns separated by nine
native indices. Its certificates carry those inner identities, certificates
and scope map; an unsupported outer tag is rejected. Exposure during outer
closure preserves both inner wave identities without making a shedding claim.
For the positive retained lane only, a declared exact toy bridge binds the
committed response event to closure load
\(L^{\rm toy}=C^{\rm cap}+J\) and closure cost
\(C_{\rm close}^{\rm toy}=C^{\rm cap}\). It does not identify current with
SADAR pressure. The resulting exact residual and surplus are both \(1\); the
computed record binds the source state/event, outer certificate, inner
lineages, exoshedding route back to their \(B0\) scope, commitment order, and
operand digest. Only that record authorizes sheddic exposure, which preserves
both inner wave identities. Replacing the fold status by the string
\texttt{shedding} is rejected. The matched removed lane has no fold and no
shedding record; exact restoration reproduces the complete record byte for
byte. The retained inner relation is then carried with the exact fold witness
and higher contact provenance.
At the higher scope, retention is committed at \(56\), the new intervention
boundary is \(58\), its changed response is at \(59\), and the next outer fold
starts at \(60\). Removing the higher retained coordinate blocks that offer;
restoring it restores the same admitted construction. A local coincident
trace fails the outer-fold and cross-scale claims. Action-origin retention
also carries its authentic source action registry and commitment evidence.

\subsection{Passive observation, interrogation and operational mindfulness}
\label{subsec:query-mindfulness}
\paragraph{Setup.}
A passive observation reads a committed state after the transition and is
not an input to \(\mathcal U\). Interrogation is instead an explicit
perturbation \(q_k\), with the observation on which it depends already
committed through \(k\). Queried and unqueried inputs are distinct
perturbation families; an observer record alone cannot change execution.

\paragraph{Derivation.}
For this specialization, the operative boundary of \mainequation{main:eq:operative-self}
is established by the matched causal criterion and remains stable over the
observation window. Consider a reflexive query \(q_k=-J_k\), with current
observed at \(k\) and response at \(k+1\). For constant forcing
\(b=A+g\sum r_e\), \eqref{eq:awareness-response} becomes
\begin{equation}
 J_{k+1}=\left(1-\frac1p\right)J_k+\frac bp,\qquad
 J_{k+1}-b=\left(1-\frac1p\right)(J_k-b).
 \label{eq:mindfulness-stability}
\end{equation}
Its constant-forcing convergence holds exactly when \(p>1/2\). At \(p=1\)
and \(b=1\), every subsequent current is exactly one while the declared H1
and hinge task continues. An extra delay using \(J_{k-1}\) instead has a
different recurrence and can oscillate. Its query retains the source cut
\(k-1\); the initial history value is declared before the first query.
Variable forcing cannot silently
inherit the constant-forcing conclusion.

\paragraph{Specialization.}
Operational mindfulness denotes this stable reflexive re-addressing of an
already operative state of this model while task dynamics continue.

\paragraph{Falsification.}
A query without a committed observation is incomplete. Stability alone
does not establish provenance or the earlier causal operativity witness.

\aodprovenance{\eqref{eq:awareness-response};
\mainequation{main:eq:causal-update};
\eqref{eq:mindfulness-stability}.}

\subsubsection{Query and mindfulness executions}
\label{manual:m2-query-mindfulness}
A passive observer runs after each committed transition. Interrogation is a
separate input \(q\) with a declared committed observation and changes the
perturbed execution. It cannot be counted as passive tracing.
The established operative boundary is stable over this observation window;
stabilization before that boundary is established is not a mindfulness witness.
For \(p=1\), constant retained forcing one and \(q_k=-J_k\), an initial
current three becomes one and remains one while H1/fibre transitions continue.
The observation is of current \(J_k\), already committed at \(k\); the response
is at \(k+1\). An extra delay using \(J_{k-1}\) and the boundary pressure
\(p=1/2\) both fail this stabilization conclusion. The general constant-forcing
criterion \(p>1/2\) is derived in \eqref{eq:mindfulness-stability}.
The delayed control declares \(J_{32}=0\) as its initial history and
\(J_{33}=3\) as its starting current. Its first query cites cut 32;
the query at cut 34 uses \(J_{33}=3\) and cites cut 33. Each event retains
the consumed observation cut. The initialized history is a model premise;
subsequent query values come from the recorded preceding observation.

\subsection{Preserved example ladder and source status}
The original A0--A18 examples share the constructors of
\secref{sec:m2-awareness-ladder} and its declared shared runner
\texttt{manual/scripts/aod\_awareness.py}. Their required
roles remain explicit in the following crosswalk.
\begin{longtable}{p{0.12\linewidth}p{0.36\linewidth}p{0.40\linewidth}}
\toprule Tier & Construction & Current witness or control\\\midrule
\endhead
A0 & Inherited closure and temporal interfaces & Exact native source binding and inherited suite\\
A1 & Passive trace & Identical state/event bytes, observation separated\\
A2 & Structural ancestry & Closure memory without retention effect\\
A3 & Matched channel operativity & Matched pressure families, removal/restoration\\
A4 & Attention separation & Present attention one, retention effect false\\
A5 & Persistence window & Surviving relation with and without later gain\\
A6 & Polycadence & Lock, nonzero bias, unequal cadence and jitter\\
A7 & Crescendo/recovery & Operative-source counts \((1,2,3,0,3)\)\\
A8 & Later contact/fold/chain & Two separate causal interventions and B2 offer\\
A9 & Ternary boundary roles & SELF/BETWEEN/NOT-SELF with distinct Null type\\
A10 & Indexical/shared labels & \emph{I} under node renaming; impostor rejection\\
A11 & Sensor--motor & Issued command and admitted channel\\
A12 & Proprioception & Bound return projection and later residual effect\\
A13 & CPU/subnet & Operative interface and separate operation-count proxy\\
A14 & Learning & Persistent \((5,6,7)\) versus removed \((4,4,4)\)\\
A15 & Reflexive interrogation & Explicit query perturbation with prior observation\\
A16 & Operational mindfulness & Exact stabilization under stated hypotheses\\
A17 & Fold/expose/carriage & Inner identities survive; higher retained pair passes\\
A18 & Recurrent other/subnet & Separate admitted peer and repeated contact prefixes\\
\bottomrule
\end{longtable}

The forty-seven new input rows have identities \texttt{R04-NEW-001} through
\texttt{R04-NEW-047}. Their source, scope, record class, phase type,
control/trace/perturbation fields, expected exact result and falsification
condition are included in the JSONL data. They are newly authored witnesses.
The separate nineteen-row historical source roster remains
\texttt{HOLD\_MISSING\_AUTHENTIC\_SOURCES}. No original script was reconstructed
from a basename or a former runtime, and these new executions do not satisfy
historical rerun rows.

\aodprovenance{\mainequation{main:eq:awareness-state-factorization};
\mainequation{main:eq:retained-rebalancing}; \mainequation{main:eq:awareness-active};
\mainequation{main:eq:awareness-operative-set};
\eqref{eq:awareness-response}; \eqref{eq:awareness-causal-difference};
\mainequation{main:eq:operative-self}; \eqref{eq:action-return-retention};
\eqref{eq:awareness-shedding-bridge};
\eqref{eq:awareness-conditioned-shedding};
\mainequation{main:eq:awareness-carriage}; \eqref{eq:mindfulness-stability};
\mainequation{main:eq:closure-phase}; \mainequation{main:eq:wave-identity}; \mainequation{main:eq:relational-time};
\mainequation{main:eq:fold-outer-closure}; \mainequation{main:eq:expose-inner-lineage}.}
